Small treewidth alone does not make nonconvex mixed-integer optimization
tractable: box-constrained nonconvex quadratic programming is strongly
NP-hard at treewidth two. We show that two further quantities make tree
decompositions useful for certified global optimization: upper bounds $L_i$
on the curvature along each coordinate, and quadratic growth. Subtracting
$L_iw^2/8$ at each grid node, where $w$ is the longest adjacent grid
interval, turns exact dynamic programming over a product grid into a valid
lower bound. Min-marginals then filter the grids, and the stage
grids form a certificate that is checked by recomputation and needs no
growth assumption. Under growth in the norm weighted by the $L_i$, graded
grids keep $O(\sqrt{\bar\kappa}\log(n+2))$ nodes per coordinate at every
accuracy, where $\bar\kappa$ is a scale-invariant ratio of curvature to
growth. For rational mixed-integer quadratic programs on a box with bag size
$p$ and input length $I$, this certifies accuracy $2^{-q}$ in
$f(p,\bar\kappa)(I+q+1)^5$ bit operations; under growth at a unique
minimizer, with the Euclidean ratio $\kappa\ge\bar\kappa$, an exact
minimizer takes $f_1(p,\kappa)I^{O(1)}$. The growth constant is never
needed. A moderate $\kappa$ requires the Hessian block of the interior continuous
coordinates to have smallest eigenvalue at least $2L/\kappa$, and distant
local minima to be clearly worse than the global one. Unless P${}={}$NP, neither parameter can be
dropped and the dependence on $\kappa$ cannot be polylogarithmic; under the
randomized exponential-time hypothesis its exponent cannot be $o(p)$. We
extend the method to exact partial minimization, totally unimodular coupling
constraints and several minimizers, and test an exact-arithmetic
implementation with an independent checker.
